% !TEX root = ../oddp.tex

\section{Homotopy cyclic diagonals and power operations}\label{s:2bdiagonals}

Let $\rho$ denote the standard generator of the cyclic group $\Cyc_r$ on $r$ elements, which we identify with $\{0,1,\dots,r-1\}$ via $i\mapsto \rho^i$.
The symmetric group on $\{0,1,\dots,r-1\}$ is denoted $\Sym_r$, and the regular representation includes $\Cyc_r$ into $\Sym_r$.
The letter $\tilde{r}$ will denote the quantity $\frac{r-1}{2}$.
The definitions of \cref{s:unstable} are taken from \cite{may1970general} and \cite{medina2021may_st}.

\subsection{Homotopy cyclic diagonals}\label{s:unstable}

The \emph{minimal $\Cyc_r$-resolution} of the commutative ring $R$ with trivial action is the chain complex $\W{r} = \bigoplus_{q \geq 0} \W{r}_q$ defined as follows:
\begin{align*}
	\W{r}_q &= R[\Cyc_r]\langle e_q\rangle &
	\partial(e_q) &= \begin{cases}
		N(e_{q-1}) & \text{if $q$ is even} \\
		T(e_{q-1}) & \text{if $q$ is odd,}
	\end{cases}
\end{align*}
where $N = \sum_j \rho^j$ and $T = \rho - \Id$.

\begin{definition}
	A \emph{homotopy cyclic $r$-diagonal} on a chain complex $\ucadenas$ is a $\Cyc_r$-equivariant chain map
	\[
	\mu_r \colon \W{r} \ot \ucadenas \lra \ucadenas^{\ot r}.
	\]
\end{definition}

Let $\ES$ be the Barratt--Eccless operad. Suppose now that $\ucadenas$ is a $\ES$-algebra. A \emph{May--Steenrod structure} is a family of homotopy cyclic $r$-diagonals $\{\mu_r\}_{r>1}$ that factor as
\[
\W{r} \ot \ucadenas \xra{g_r \ot \id} \ES(r) \ot \ucadenas \lra \ucadenas^{\ot r}
\]
$g_r$ is equivariant with respect to the regular representation of $\Cyc_r$, and the second map is the structural map of the $\ES$-algebra.	

\begin{proposition}\cite{may1970general}\label{prop:unstable}
	If $r$ is prime and $R=\bF_r$, a homotopy cyclic $r$-diagonal $\mu_r$ on a chain complex $\ucadenas$ induces operations indexed by the integers
	\begin{align*}
		\power^{i} \colon H_{-m}(\ucocadenas;\mathbb{F}_r)& \lra H_{-m-i}(\ucocadenas;\mathbb{F}_r)
	\end{align*}
	that send a cohomology class $[x]$ in degree $-m$ to the cohomology class $[y]$ with
	\[
	y(a) = (-1)^{\binom{m}{2} \tilde{r}}(\tilde{r}!)^m(x \ot \overset{r}{\dots} \ot x)(\mu_r(e_{(r-1)m-i} \ot a)).
	\]
	If additionally the homotopy cyclic $r$-diagonal is part of a May--Steenrod structure, then these operations satisfy the following
	\begin{itemize}
		\item $\power^i(x) = 0$ if $i>(r-1)m$ and $x$ has degree $-m$.
		\item $\power^i = 0$ unless $i = 2(r-1)k$ or $i = 2(r-1)k+1$.
	\end{itemize}
\end{proposition}
%
%\begin{proof}
%	See . For the last claim, note that if $i>(r-1)m$, then $(r-1)m-i$ is negative, and therefore $\mu(e_{(r-1)m-i},a)$ vanishes.
%\end{proof}

\begin{example}
	The cup-$i$ coproducts of Steenrod \cite{steenrod1947products, medina2022axiomatic} yield a natural May--Steenrod structure of arity $2$ on the normalized chains of simplicial sets which is May--Steenrod.
	For general $r$, a similar May--Steenrod unstable $r$-diagonal was constructed in \cite{medina2021may_st}.
\end{example}

\subsection{Stable homotopy cyclic diagonals}\label{s:stable}

Let $\Wst(r)$ be the unbounded chain complex
\[
\dotsb \overset{N}{\lla} R[\Cyc_r]\langle e_{-2}\rangle
\overset{T}{\lla}
R[\Cyc_r]\langle e_{-1}\rangle
\overset{N}{\lla}
R[\Cyc_r]\langle e_0\rangle
\overset{T}{\lla}
R[\Cyc_r]\langle e_{1}\rangle
\overset{N}{\lla}
\dots
%R[\Cyc_r]\langle e_{2}\rangle
%\overset{T}{\lla}
\]
If either $r$ is odd or $r=2$ and $R=\bF_2$, we understand it as the limit of
\[
\dotsb \xra{\eta} \sus{2(1-r)}\W{r} \xra{\eta} \sus{1-r} \W{r} \xra{\eta} \W{r}.
\]
%In general, we understand it as the limit of
%\[
%\dotsb \lra \sus{4(1-r)}\W{r}\overset{\eta_{2(1-r)}}{\lra} \sus{2(1-r)} \W{r} \overset{\eta_{2(1-r)}}{\lra} \W{r}.
%\]
where $\eta \colon \W{r} \to \sus{r-1}\W{r}$ is the \emph{suspension homomorphism} that sends $e_{q}$ to $e_{q+1-r}$. 

\begin{definition}
	A \emph{stable homotopy cyclic $r$-diagonal} on a chain complex $\ucadenas$ is a $\Cyc_r$-equivariant chain map
	\[
	\mu_r \colon \Wst(r) \ot \ucadenas \lra \ucadenas^{\ot r}.
	\]
\end{definition}

	Let $\ESst_r$ be the inverse limit of
	\[
	\dotsb \lra \susp{2(1-r)} \ES_r\overset{\eta}{\lra}\susp{1-r} \ES_r\overset{\eta}{\lra} \ES_r
	\]
	where $\eta$ is the operadic suspension of \cite{berger2004combinatorial} that sends a sequence of permutations $(\sigma_0,\dots,\sigma_q)$ to zero if $(\sigma_0(0),\dots,\sigma_{r-1}(0))$ has repeated elements, and to $\pm (\sigma_{r-1},\sigma_r,\dots,\sigma_q)$ otherwise, where $\pm$ is the sign of the permutation that orders the sequence $(\sigma_0(0),\dots,\sigma_{r-1}(0))$.	Suppose now that $\ucadenas$ is an algebra over $\ESst$. A $\emph{stable May--Steenrod}$ structure is a family of stable homotopy cyclic $r$-diagonals $\mu_r\colon \Wst(r)\ot \ucadenas\to \ucadenas^{\ot r}$ that factor as
	\[
	\Wst(r) \ot \ucadenas \xra{g_r \ot \id} \ESst_r \ot \ucadenas \lra \ucadenas^{\ot r}
	\]
	where each $g_r$ is equivariant with respect to the inclusion $\Cyc_r \subset \Sym_r$ and the second map the structural map of the $\ESst$-algebra.






\begin{proposition}
	If $r$ is prime, a stable cyclic $r$-diagonal $\mu_r$ on a non-negatively graded chain complex $\ucadenas$ induces operations indexed by the integers
	\begin{align*}
		\power^{i} \colon H_{-n}(\ucocadenas;\mathbb{F}_r)& \lra H_{-n-i}(\ucocadenas;\mathbb{F}_r)
	\end{align*}
	that send a cohomology class $[x]$ in degree $-m$ to the cohomology class $[y]$ with
	\[
	y(a) = (-1)^{\binom{m}{2} \tilde{r}}(\tilde{r}!)^m(x \ot \overset{r}{\dots} \ot x)(\mu(e_{(r-1)m-i} \ot a))
	\]
\end{proposition}

\begin{proof}
	If $\ucadenas$ has a stable diagonal, then $\susp{} \ucadenas$ has a stable diagonal as well, defined as $\mu(e_q \ot \susp{} a) = \susp{\ot r}\mu(e_{q+r-1} \ot a)$.
	Therefore, after enough suspensions, we may assume that the stable diagonal vanishes in negative degrees and that $q$ is positive.
	Then a truncation reduces this to the situation of an unstable diagonal.
\end{proof}

\begin{example}[Suspension spectra]
	Every May--Steenrod structure yields a stable May--Steenrod structure by precomposing with the projections $\Wst(r) \to \W{r}$, $\ESst\to \ES$. These May--Steenrod structures vanish on $\Wst[i](r) \ot \ucadenas$ for $i<0$ (See \cref{figure:1}). The chain complex of the suspension spectrum of a pointed simplicial set $X$ is isomorphic to the chain complex of $X$, and the May--Steenrod structure on the latter gives a stable May--Steenrod structure on the former. 
\end{example}

\begin{example}[$\Sigma$-spectra]
	\cite{Gill2020} endowed the chain complex of a $\Sigma$-spectrum of pointed simplicial sets with the structure of a $\ESst$-coalgebra. It is easy to see that the May--Steenrod structure of \cite{medina2021may_st} gives also a stable May--Steenrod structure in this case.
\end{example}



\begin{remark}[Tate diagonal]\federico{Needs review}
Let $p$ be a prime number. The norm map $\nu \colon W(p) \xrightarrow{\sim} W^\dd(p)$ is the chain map that sends $e_0$ to $N \cdot e_0^\dd$ and is $0$ in both positive and negative degrees. Its mapping cone, which we denote by $\widehat W(p)$, is the \emph{Tate complex} of $\Cyc_p$. If $M$ is a $R[\Cyc_p]$-module, the \emph{Tate cohomology groups} of $\Cyc_p$ with coefficients in $M$ are computed as 
\[
\widehat H^*(\Cyc_p;M)
\cong 
H_*(\widehat W(r)\otimes_{R[\Cyc(p)]}M).
\]
There is an isomorphism of chain complexes $(\Wst)^\dd(p)\cong \widehat W (p)$. Since $\Wst(p)$ is projective, any equivariant stable comultiplication $\Wst(p)\otimes \ucadenas\to\ucadenas^{\otimes p}$ gives rise to an equivariant chain map $\ucadenas\to \widehat W(p)\otimes \ucadenas^{\otimes p}$, and quotienting by the group action yields a map 
\begin{equation}\label{eq:tate1}
	C\lra \left(\widehat W(p)\otimes \ucadenas^{\otimes p}\right)_{\h\Cyc_p}.
\end{equation}
Now, if $X$ is a spectrum, the Tate construction $\left(X^{\wedge p}\right)^{t\Cyc_p}$ on $X$ is defined as the cofiber of the norm map 
\[
(X\wedge \overset{p}{\dots} \wedge X)_{h\Cyc_p}\to (X\wedge \overset{p}{\dots} \wedge X)^{h\Cyc_p}.
\]
In \cite{Nikolaus-Scholze} it was proven that this construction is compatible with smash products and that there exist a unique natural transformation $\Delta_X\colon X\to \left(X^{\wedge p}\right)^{t\Cyc_p}$ called the \emph{Tate diagonal}. We can view the chain complex $C$ of a spectrum as the $HR$-module $X\wedge HR$. Then the smash product of the Tate diagonals
\[
	X\wedge HR
	\lra
	\left(X^{\wedge p}\right)^{t\Cyc_p}
	\wedge
	\left(HR^{\wedge p}\right)^{t\Cyc_p}
	\simeq
	\left(X^{\wedge p}
	\wedge
	HR^{\wedge p}\right)^{t\Cyc_p}
\]
composed with the map that arranges the factors
\[
\left(X^{\wedge p}
\wedge
HR^{\wedge p}\right)^{t\Cyc_p}
\lra
\left((X\wedge HR)^{\wedge p}\right)^{t\Cyc_p}
\]
and with the canonical map
\[
\left((X\wedge HR)^{\wedge p}\right)^{t\Cyc_p}
\lra
\left((X\wedge HR)^{\ot p}\right)^{t\Cyc_p},
\]
(where $\ot$ denotes the tensor product of $HR$-modules) gives a map
\begin{equation}\label{eq:tate2}
	X\wedge HR \lra \left(\left(X\wedge HR\right)^{\ot p}\right)^{t\Cyc_p}.
\end{equation}
Viewing the chain complex $C$ as a model of $X\wedge HR$, the map 
\eqref{eq:tate2} is an instance of a map of the form \eqref{eq:tate1} and therefore of a stable homotopy cyclic $p$-diagonal.
\end{remark}

\begin{figure}[b]
	\input{tables/unstable}
	\caption{A depiction of $\Wst(3) \ot \ucadenas$. A stable May--Steenrod structure comes from a May--Steenrod structure if it vanishes on the part coloured in green.}
	\label{figure:1}
\end{figure}

\begin{figure}
	\input{tables/hat}
	\caption{If a stable homotopy $r$-cyclic comultiplication for $r=3$ vanishes on the part coloured in green, then $\power^i$ vanishes for $i<0$.}
	\label{figure:2}
	
	\vspace{1cm}
	
	
\[
	\xymatrix@R=.4cm@C=.4cm{
		\phantom{W_4 \ot C_{-3}} &
		\phantom{W_4 \ot C_{-3}} &
		\phantom{W_4 \ot C_{0}} &
		\phantom{W_4 \ot C_{3}} &
		\bF_3 \ot C_{6}\ar[d]
		\\
		\phantom{W_3 \ot C_{-3}} &
		\phantom{W_3 \ot C_{-3}} &
		\phantom{W_3 \ot C_{0}} &
		\phantom{W_3 \ot C_{3}} &
		\rWd[-1] \ot C_{6}\ar[d]
		\\
		\phantom{W_2 \ot C_{-3}} &
		\phantom{W_2 \ot C_{-3}} &
		\phantom{W_2 \ot C_{0}} &
		\bF_3 \ot C_{3}\ar[d] &
		\rWd[-2] \ot C_{6}\ar[d]\ar[l]
		\\
		\phantom{W_1 \ot C_{-3}} &
		\phantom{W_1 \ot C_{-3}}&
		\phantom{W_1 \ot C_{0}} &
		\rWd[-1] \ot C_{3}\ar[d] &
		\rWd[-3] \ot C_{6}\ar[d]\ar[l]
		\\
		\phantom{W_0 \ot C_{-3}} &
		\phantom{W_0 \ot C_{-3}}&
		\bF_3 \ot C_{0}\ar[d] &
		\rWd[-2] \ot C_{3}\ar[d]\ar[l] &
		\rWd[-4] \ot C_{6}\ar[d]\ar[l]
		\\
		\phantom{\rWd[-1] \ot C_{-3}} &
		\phantom{\rWd[-1] \ot C_{-3}} &
		\rWd[-1] \ot C_{0}\ar[d] &
		\rWd[-3] \ot C_{3}\ar[d]\ar[l] &
		\rWd[-5] \ot C_{6}\ar[d]\ar[l]
		\\
		\phantom{\rWd[-2] \ot C_{-3}} &
		\bF_3 \ot C_{-3} &
		\rWd[-2] \ot C_{0}\ar[l] &
		\rWd[-4] \ot C_{3}\ar[l] &
		\rWd[-6] \ot C_{6}\ar[l]
}
\]

	\caption{The complex $\susp{\bs(r-1)}\Wdualaug{r} \ot \ucadenas$ for $r=3$ and $R=\bF_3$. The map $\varphi$ goes from Figure \ref{figure:2} to here, is an epimorphism and its kernel is the smallest chain subcomplex that contains the green part. The power operation $\power^0$ is computed at the corners $\bF_3\ot C^3_n$. Due to the presence of suspensions, note that the bidegree of $e_q^\dd\ot x_n$ is $ (n(r-1) - q, n)$.} 
	\label{figure:3}
\end{figure}

\subsection{Connected homotopy cyclic diagonals}\label{s:connected}

Let $\mu_r$ be a stable homotopy cyclic $r$-diagonal on $\ucadenas$.
A way of guaranteeing the desirable property that $\power^i(x)$ vanishes for $i<0$ is to ask for $\mu_r$ to vanish on
	\[
	\bigoplus_n\bigoplus_{q>(r-1)n}\Wst[q](r) \ot \ucadenas[n].
	\]
We will now identify the quotient of $\Wst(r)\otimes \ucadenas$ by this graded submodule when $r$ is odd or $r=2$ and $R=\bF_2$ (see Figures \ref{figure:2}). We then define a connected homotopy cyclic $r$-diagonal as a map from this quotient to $\ucadenas^{\otimes r}$.



Let $\Waug{r}$ be the augmented minimal resolution of $\Cyc_r$, i.e., the chain complex
\[
R
\overset{\varepsilon}{\lla} 
R[\Cyc_r]\langle e_{1}\rangle
\overset{T}{\lla}
R[\Cyc_r]\langle e_{2}\rangle
\overset{N}{\lla}
R[\Cyc_r]\langle e_{3}\rangle
\overset{T}{\lla}
\dotsb
\]
(note the reindexing, so that $e_i$ has degree $i$) and $\Wdualaug{r}$ for its dual:
\[
\dotsb
\overset{T^\dd}{\lla}
R[\Cyc_r]\langle e_{3}^\dd\rangle
\overset{-N}{\lla}
R[\Cyc_r]\langle e_{2}^\dd\rangle
\overset{T^\dd}{\lla}
R[\Cyc_r]\langle e_{1}^\dd\rangle
\overset{-\varepsilon}{\lla} 
R,
\]
where $T^\dd(e_i^\dd)=\left((\rho^{-1}-1)\cdot e_{i+1}\right)^\dd$.
Observe that $\rho$ acts with the dual action in the dual complex, so $\rho \cdot e_{i}^\dd = (\rho^{-1} e_i)^{\dd}$.
The generator of the copy of $R$ in degree $0$ will be denoted as $1$ or as $e_0^\dd$.
If $r$ is odd or $r=2$ and $R=\bF_2$, there is a chain map $\theta \colon \Waug{r} \to \sus{1-r}\Waug{r}$ defined as 
\begin{align*}
	\theta(e_0) &= N(e_{r-1}),
	&
	\theta(e_{i}) &= e_{i+r-1}, \quad i>0 
\end{align*}
whose dual $\theta^\vee\colon \Wdualaug{r}\to \Sigma^{1-r}\Wdualaug{r}$ takes the form
\begin{align*}
	\theta(e_{i}^{\vee}) &= 0, \quad i < r-1
	&
	\theta(e_{i}) &= e_{i+r-1}, \quad i\geq r-1.
\end{align*}

Let $(I,>)$ be the poset of natural numbers (i.e., non-negative integers). A \emph{directed system of $R$-modules} is a functor $F_\bs\colon I\to \Mod{R}$. If $F_\bs,G_{\bs}$ are two directed systems, then $F_\bs\ot G_{\bs}$ denotes the directed system obtained by performing the tensor product levelwise. Maps between directed systems are natural transformations of functors. 

A chain complex $C$ can be viewed as a directed system of $R$-modules $C_{\bs}$ setting $C_i$ as the $R$-module of chains of degree $i$, and letting the differential be the structure map. 

From this viewpoint, chain complexes are those directed systems for which the composition of any two structure maps is zero. Observe that the tensor product of a directed system and a chain complex is again a chain complex.

Similarly, if $F_\bs$ is a directed system of chain complexes and $G_\bs$ is a directed system of $R$-modules, then $F_\bs \otimes G_\bs$ is a bicomplex.

\begin{definition} Let $\susp{\bs(r-1)}\Wdualaug{r}$ be the directed system of chain complexes
\[
	\Wdualaug{r}\overset{\theta^\vee}{\lla }
	\susp{r-1}\Wdualaug{r}\overset{\theta^\vee}{\lla }
	\susp{2(r-1)}\Wdualaug{r} \overset{\theta^\vee}{\lla }
	\ldots
	\overset{\theta^\vee}{\lla }	
	\susp{n(r-1)}\Wdualaug{r} \overset{\theta^\vee}{\lla }
	\ldots
\]
\end{definition}

\begin{example}\label{ex:311}
	Let $C$ be a $R$-chain complex, which we view as a $R[\Cyc_r]$-chain complex with trivial $\Cyc_r$ action.  The tensor product $\susp{\bs(r-1)}\Wdualaug{r} \ot C_\bs$ is bicomplex of $R[\Cyc_r]$-modules whose totalization has the following differential:
	\[
	\partial(e^\dd_{q} \ot \vec{x}) = \partial (e^\dd_{q}) \ot \vec{x} + (-1)^q \theta^\dd(e^\dd_{q}) \ot \partial (\vec{x}).
	\]
\end{example}
There is a functor $Q\colon \asimplexinj\to I$ that sends an ordinal $[m]$ to the number $m$. Thus, any ordered sequence of chain complexes $F$ yields, precomposing with this functor, a semi-simplicial chain complex $Q^*(F)$, and we have:
\begin{lemma} \label{lemma:adj}
	If $X_\bullet$ is a semi-simplicial $R$-module, then there is a canonical isomorphism
\begin{equation}\label{eq:adj}
	\chains(Q^*(F)\ot X_\bullet^{\nd})\cong F\ot \chains(X). 
\end{equation}
\end{lemma}

%If $r$ is even and $R$ is arbitrary, define the cochain complex $\rWhatd(r)$ as
%\[
%\dotsb
%\overset{-T}{\lla} 
%R[\Cyc_r]\langle e_{-2}^\dd\rangle 
%\overset{N}{\lla}
%R[\Cyc_r]\langle e_{-1}^\dd\rangle
%\overset{-T}{\lla} 
%R[\Cyc_r]\langle e_{0}^\dd\rangle / N 
%\]
%This chain complex admits chain maps
%\begin{align*}
	%\hat{\theta}_{r-1} \colon \rWd(r)& \lra \susp{1-r}\rWhat(r)
	%&
	%\hat{\theta}_{r-1} \colon \rWhat(r)& \lra \susp{1-r}\rW(r).
%\end{align*} 	%%They have the same name?
%Define $\ucadenas[\mathrm{even}] = \bigoplus_{m}\ucadenas[2m]$ and $\ucadenas[\mathrm{odd}] = \bigoplus_{m}\ucadenas[2m+1]$ and write $\rWd(r) \hotimes \rucadenas$ for the graded $R[\Cyc_r]$-module $\left(\rWd(r) \ot \rucadenas[\mathrm{odd}]\right)\oplus \left(\rWhatd(r) \ot \rucadenas[\mathrm{even}]\right)$ with differential
%\[
%\partial(e^\vee_{-q} \ot \vec{x}) = \partial (e^\vee_{-q}) \ot \vec{x} + (-1)^q \theta_{1-r}(e^\vee_{-q}) \ot \partial (\vec{x}).
%\] 	%%Shouldn't it be $\hotimes_{1-r}...$?
%Observe that if $r=2$ and $R=\bF_2$, the complexes $\rWd(r)$ and $\rWhatd(r)$ coincide, and so do both definitions of $\rWd(r) \hotimes \rucadenas$.

Define the chain map $\varphi \colon \Wst(r)\ot \ucadenas \to \susp{\bs(r-1)}\Wdualaug{r} \ot C$ as 
\[
\varphi(e_q \ot x) = (-1)^{\binom{n(r-1)-q+1}{2}}e_{q-(r-1)n}^\dd \ot \vec{x}
\]
if $x$ has degree $n$.\federico{The sign needs revision.} Observe that
\begin{lemma} The map $\varphi$ is an epimorphism with kernel the smallest chain subcomplex that contains $\bigoplus_n\bigoplus_{i>(r-1)n} \Wst[i] \ot \ucadenas[n]$.
\end{lemma}


\begin{definition}\label{def:connected_diagonal}
If $r$ is odd or $r=2$ and $R=\bF_2$, a \emph{connected homotopy cyclic $r$-diagonal} on a chain complex $\ucadenas$ is a $\Cyc_r$-equivariant chain map
	\[
	\mu \colon \susp{\bs(r-1)}\Wdualaug{r} \ot C_{\bs} \lra \ucadenas^{\ot r}.
	\]
\end{definition}



\federico{He comentado una explicación no concluyente de cómo se obtiene esta aplicación. No sé si merece la pena liarse con esto}
%If either $k$ is even or $r=2$ and $R=\bF_2$, let $\varphi_k \colon \Wst(r) \to \susp{-k}\rWd(r)$ be the chain map that sends $e_q$ to $(-1)^{\binom{k-q}{2}}e_{k-q}^{\vee}$.
%One can obtain this chain map as follows: We already saw $\Wst(r)$ as the homotopy fibre of the norm map $N \colon \W{r} \to \Wdual{r}$, but it is also the homotopy fibre of the map $\Wdualaug{r} \to \Waug{r}$ that is the identity on $R$ and zero elsewhere.
%The map $\varphi_k$ is the composition
%\[
%\Wst(r)\overset{\theta_k}{\lra} \sus{-k} \Wst(r) \overset{\cong}{\lra} \sus{-k}\hofib(\rWd(r) \to \rW(r)) \lra \susp{-k}\rWd(r).
%\]	%%Why $\sus$ instead of $\susp$? Aren't the same for $k$ even? $\theta_{k}$ is defined with $\susp$.
%This composition is an epimorphism with kernel the smallest chain subcomplex that contains $\Wst[>k]$.
%In other words, it is isomorphic to the truncation $\tau_k \Wst(r)$.
%%If $k$ is odd, let $\varphi_k \colon \Wst(r) \to \sus{k}\rWhatd(r)$ be the chain map that sends $e_q$ to $e_{k-q}^{\vee}$.

As a consequence,

\begin{proposition}
	If $r$ is prime and $R=\bF_r$, a connected homotopy cyclic $r$-diagonal induces operations indexed by the integers
	\begin{align*}
		\power^{i} \colon H_{-m}(\ucocadenas;\mathbb{F}_r)& \lra H_{-m-i}(\ucocadenas;\mathbb{F}_r)
	\end{align*}
	that send a cohomology class $[x]$ in degree $m$ to the cohomology class $[y]$ with
	% (-1)^{\binom{ir}{2} + \binom{m}{2}\tilde{r}}(\tilde{r}!)^m
	\[
	y(a) = (x \ot \overset{r}{\dots} \ot x)(\mu(e_{ri}^\vee \ot a))
	\]
	These operations satisfy the following equation
	\begin{itemize}
		\item $\power^i(x) = 0$ for $i<0$.
	\end{itemize}
\end{proposition}

\begin{example}
	As far as we know, the concept of connected homotopy cyclic $r$-diagonal has not appeared before in the literature, nonetheless, this viewpoint is implicit in the construction of the homotopy comultiplication for $r=2$ on the cochain operations on the normalized cochains of a simplicial set of \cite{medina2021fast_sq}.
	It is also implicit in the construction of the stable homotopy comultiplication on the chain complex of augmented semi-simplicial objects in the Burnside $2$-category \cite{cantero-moran2020khovanov}.
\end{example}

\begin{question}
	Give a good definition of ``May--Steenrod connected $r$-diagonal''.
	Give a homotopy equivalence between the Tate complex of $\ES_r$ and the unbounded complex $\ESst_r$.\federico{The second one, I find it challenging, but unmotivated in the current version of the paper.}
\end{question}

